\section*{Hoofdstuk \ref{ch:b1:complexation} --- Complexvorming}

\begin{solution}{exo:b1:complexation:1}
$[\ce{Cu(NH3)4}]^{2+}$: \ce{Cu^2+}, vier \ce{NH3} gebonden via N.
$[\ce{Fe(CN)6}]^{4-}$: \ce{Fe^2+}, zes \ce{CN-} gebonden via C;
$[\ce{Fe(CN)6}]^{3-}$: hetzelfde met \ce{Fe^3+}. $[\ce{Ag(NH3)2}]^{+}$:
\ce{Ag+}, twee \ce{NH3} via N. $[\ce{FeSCN}]^{2+}$: \ce{Fe^3+}, één
\ce{SCN-}, gebonden via N, ondanks de manier waarop de formule geschreven
wordt. $[\ce{Zn(OH)4}]^{2-}$: \ce{Zn^2+}, vier \ce{OH-} via O.
$[\ce{CaY}]^{2-}$: \ce{Ca^2+}, één edta-ion via twee N en vier O.
\end{solution}

\begin{solution}{exo:b1:complexation:2}
$\log K_i = \log\beta_i - \log\beta_{i-1}$: 4.10, 3.41, 2.82, 2.03, en
$\mathrm{p}K_{d,i} = \log K_i$ neemt dezelfde waarden aan. De laatste
heeft betrekking op
$[\ce{Cu(NH3)4}]^{2+} \rightleftharpoons [\ce{Cu(NH3)3}]^{2+} + \ce{NH3}$.
\end{solution}

\begin{solution}{exo:b1:complexation:3}
Gebieden: $[\ce{Cu(NH3)4}]^{2+}$ onder pNH3 2.03, dan $n = 3$ tot 2.82,
$n = 2$ tot 3.41, $n = 1$ tot 4.10, \ce{Cu^2+} erboven. pNH3 = 0:
$[\ce{Cu(NH3)4}]^{2+}$; 2.5: $[\ce{Cu(NH3)3}]^{2+}$; 3.5:
$[\ce{CuNH3}]^{2+}$; 5.0: \ce{Cu^2+}.
\end{solution}

\begin{solution}{exo:b1:complexation:4}
$[\ce{FeSCN^2+}]/[\ce{Fe^3+}] = \beta_1[\ce{SCN-}] = 10^{2.96} \times 0.10 =
91$: de gecomplexeerde fractie is $91/92 = \qty{98.9}{\percent}$.
\end{solution}

\begin{solution}{exo:b1:complexation:5}
Elke $[\mathrm{ML}_k] = \beta_k[\mathrm{M}][\mathrm{L}]^k$; delen door de
som over $k$ (het totale metaal) elimineert $[\mathrm{M}]$. Met drie
deeltjes is $f_i = 1/(1 + [\mathrm{ML}_{i-1}]/[\mathrm{ML}_i] +
[\mathrm{ML}_{i+1}]/[\mathrm{ML}_i]) = 1/(1 + 1/(K_i x) + K_{i+1}x)$ met
$x = [\mathrm{L}]$. De noemer is het kleinst wanneer zijn afgeleide
$-1/(K_i x^2) + K_{i+1}$ nul wordt, $x^2 = 1/(K_iK_{i+1})$, waar de twee
buren gelijk zijn, en $\mathrm{pL} = \frac12(\log K_i + \log K_{i+1})$.
Voor $[\ce{Cu(NH3)2}]^{2+}$: $\frac12(3.41 + 2.82) = 3.12$.
\end{solution}

\begin{solution}{exo:b1:complexation:6}
$[\ce{FeSCN}]^{2+} + \ce{F-} \rightleftharpoons [\ce{FeF}]^{2+} + \ce{SCN-}$,
$K = 10^{6.85 - 2.96} = 10^{3.89}$. Verhouding $= K[\ce{F-}]/[\ce{SCN-}] =
10^{3.89} \times 0.010/0.10 = 780$: het rode \omterm{def:b1:complexation:complex}{complex} is bijna volledig
vervangen en de kleur verbleekt.
\end{solution}

\begin{solution}{exo:b1:complexation:7}
$[\ce{MgY}]^{2-} + \ce{Ca^2+} \rightleftharpoons [\ce{CaY}]^{2-} + \ce{Mg^2+}$,
$K = 10^{12.69 - 10.90} = 10^{1.79} = 62$. Met gelijke hoeveelheden is
$x^2/(1 - x)^2 = 62$, $x/(1 - x) = 7.9$, $x = 0.89$: \qty{89}{\percent}
van het magnesium komt vrij.
\end{solution}

\begin{solution}{exo:b1:complexation:8}
\ce{Ag2O(s) + H2O + 4NH3 <=> 2[Ag(NH3)2]+ + 2OH-}, $K = \beta_2^2 \times
10^{-15.43} = 10^{14.44 - 15.43} = 10^{-0.99}$. De constante is niet
klein: zodra de vrije ammoniak enkele tienden mol per liter bereikt, lost
het bruine oxide dat door de eerste druppels gevormd werd (waarbij de
ammoniak als base optrad) op tot het zilverammine.
\end{solution}

\begin{solution}{exo:b1:complexation:9}
pH 10: $\alpha = 1 + 10^{1.24} + 10^{-1.96} = 18.4$, $\log\beta' = 12.69 -
1.26 = 11.43$. pH 7: $\alpha = 1 + 10^{4.24} + 10^{4.04} + 10^{0.19} +
\cdots = 10^{4.45}$, $\log\beta' = 8.24$. Bij pH 7 ligt de conditionele
constante onder $10^{10}$: de reactie van calcium met edta is niet
volledig genoeg voor een titratie, omdat het meeste edta geprotoneerd is.
\end{solution}

\begin{solution}{exo:b1:complexation:10}
1.~\ce{CuO(s) + H2O + 4NH3 <=> [Cu(NH3)4]^2+ + 2OH-}, $K = \beta_4 \times
10^{-20.64} = 10^{12.36 - 20.64} = 10^{-8.28}$.
2.~$\mathrm{pH} = 7 + \frac12(9.25 + \log 1.0) = 11.63$, $[\ce{OH-}] =
10^{-2.37}$, en $[\ce{Cu(NH3)4^2+}] = 10^{-8.28}/10^{-4.75} =
\qty{3.0e-4}{mol/L}$: er lost heel weinig op.
3.~$\mathrm{pH} = 9.25$, $[\ce{OH-}] = 10^{-4.75}$, en de formule geeft
$10^{-8.28 + 9.50} = 10^{1.22}$, veel meer dan de ammoniak kan dragen: het
oxide lost op tot de ammoniak op is, en niet het evenwicht de grens
bepaalt. De ammoniumionen verbruiken het hydroxide dat bij het oplossen
vrijkomt; daarom helpen ammoniumzouten ammoniak om koperverbindingen op
te lossen.
\end{solution}

\begin{solution}{exo:b1:complexation:11}
1.~De reactie is de tweede vorming min de eerste: $K = K_2/K_1 =
10^{3.90 - 3.32} = 10^{0.58} = 3.8 > 1$, dus $[\ce{AgNH3}]^{+}$ reageert
met zichzelf wanneer het het hoofddeeltje is.
2.~Volgens \cref{exo:b1:complexation:5} ligt de grootste fractie bij
$\mathrm{pNH_3} = \frac12(3.32 + 3.90) = 3.61$, waar $\beta_1 x = 10^{-0.29}
= 0.51$ en $\beta_2x^2 = 1.0$: $f = 0.51/(1 + 0.51 + 1.0) = 0.20$. Nooit
is meer dan \qty{20}{\percent} van het zilver $[\ce{AgNH3}]^{+}$.
\end{solution}

\begin{solution}{exo:b1:complexation:12}
1.~\ce{CaCO3(s) + Y^4- <=> [CaY]^2- + CO3^2-}, $K = \beta K_s = 10^{12.69 -
8.30} = 10^{4.39}$.
2.~$x^2/(0.10 - x) = 10^{4.39}$ geeft $x = \qty{0.10}{mol/L}$ op
$4 \times 10^{-7}$ na: al het edta wordt gebruikt, \qty{10.0}{g}
calciumcarbonaat per liter.
3.~In zuur wordt edta geprotoneerd ($\mathrm{p}K_{a4} = 11.24$) en daalt
zijn conditionele constante; het carbonaat wordt ook geprotoneerd, wat
helpt, maar het is het \omterm{def:b1:complexation:complex}{complex} dat het calcium vasthoudt. Een hoge pH
houdt het edta in de vorm \ce{Y^4-}.
\end{solution}

\begin{solution}{pb:b1:complexation:1}
\textbf{1.} \ce{Ag+ + NH3 <=> [AgNH3]+}, $\beta_1 = [\ce{AgNH3+}]/([\ce{Ag+}][\ce{NH3}])$;
\ce{Ag+ + 2NH3 <=> [Ag(NH3)2]+}, $\beta_2 = [\ce{Ag(NH3)2+}]/([\ce{Ag+}][\ce{NH3}]^2)$.
\textbf{2.} $\log K_1 = 3.32$, $\log K_2 = 7.22 - 3.32 = 3.90$.
\textbf{3.} $[\ce{AgNH3}]^{+}$ zou overheersen tussen pNH3 3.90 en 3.32,
een interval dat achterstevoren loopt: het heeft geen gebied.
\textbf{4.} \ce{2[AgNH3]+ <=> Ag+ + [Ag(NH3)2]+}, $K = K_2/K_1 = 10^{0.58} =
3.8$.
\textbf{5.} $[\ce{Ag(NH3)2}]^{+}$ daaronder en \ce{Ag+} boven
$\frac12\log\beta_2 = 3.61$.
\textbf{6.} $\beta_2[\ce{NH3}]^2 = 10^{7.22 - 2.00} = 10^{5.22} = \num{1.7e5}$.
\textbf{7.} $\mathrm{p}K_d = \log\beta_2 = 7.22$.
\textbf{8.} \ce{AgCl(s) + 2NH3 <=> [Ag(NH3)2]+ + Cl-}, $K = \beta_2K_s =
10^{7.22 - 9.75} = 10^{-2.53}$.
\textbf{9.} $s = \sqrt{K}\,[\ce{NH3}] = 10^{-1.265} \times 1.0 =
\qty{0.054}{mol/L}$.
\textbf{10.} In water is $s = 10^{-4.875} = \qty{1.3e-5}{mol/L}$: ongeveer
4000 keer meer.
\textbf{11.} $0.0543 \times 143.4 = \qty{7.8}{g}$.
\textbf{12.} $[\ce{Ag+}] = K_s/[\ce{Cl-}] = 10^{-9.75}/0.054 =
\qty{3.3e-9}{mol/L} \ll s$.
\textbf{13.} $s = \sqrt{K}(1.0 - 2s)$, $s = 0.0543/(1 + 2 \times 0.0543) =
\qty{0.049}{mol/L}$.
\textbf{14.} Ammoniak van \qty{1.0}{mol/L} heeft pH 11.6, meer dan
$9.25 + 1$: maar \qty{0.4}{\percent} ervan is \ce{NH4+}.
\textbf{15.} $K = 10^{7.22 - 12.27} = 10^{-5.05}$.
\textbf{16.} $s = 10^{-2.525} = \qty{3.0e-3}{mol/L}$, \qty{0.56}{g/L}.
\textbf{17.} $K = 10^{-8.85}$, $s = 10^{-4.425} = \qty{3.8e-5}{mol/L}$,
\qty{8.8}{mg/L}.
\textbf{18.} Behandel elk neerslag met verdunde ammoniak: alleen het
chloride lost op. Geconcentreerde ammoniak lost het bromide op, trager en
gedeeltelijk; het jodide lost niet op.
\textbf{19.} $s = 1.0/187.8 = \qty{5.3e-3}{mol/L}$; $[\ce{NH3}] = s/\sqrt{K}
= \num{5.3e-3}/10^{-2.525} = \qty{1.8}{mol/L}$.
\textbf{20.} $s = \qty{4.3e-3}{mol/L}$ en $[\ce{NH3}] = \num{4.3e-3}/10^{-4.425}
= \qty{113}{mol/L}$, twee keer de concentratie van water in zuiver water
(\qty{55}{mol/L}): onmogelijk. Ammoniak kan zilverjodide niet oplossen.
\textbf{21.} \ce{[Ag(NH3)2]+ + Cl- + 2H3O+ -> AgCl(s) + 2NH4+ + 2H2O},
$K = 10^{2.53} \times (10^{9.25})^2 = 10^{21.03}$: het witte neerslag
verschijnt opnieuw, wat zilverchloride bevestigt.
\textbf{22.} $s = 1.0/143.4 = \qty{6.97e-3}{mol/L}$; $[\ce{NH3}] =
s/\sqrt{K} = \num{6.97e-3}/10^{-1.265} = \qty{0.128}{mol/L}$.
\textbf{23.} $2s = \qty{0.0139}{mol/L}$.
\textbf{24.} $[\ce{Ag+}] = 10^{-9.75}/\num{6.97e-3} = \qty{2.6e-8}{mol/L}$,
verwaarloosbaar naast $s$.
\textbf{25.} $0.128 + 0.014 = $ \textbf{\qty{0.142}{mol/L} ammoniak}:
\qty{0.142}{mol} in de liter lost precies \qty{1.0}{g} zilverchloride
op.
\end{solution}
